% =====================================================================
%  radcluster_2_1_RIS_plan.tex
%  Development plan for RadCluster_2_1: radiation-induced segregation (RIS)
%  and solute precipitation as a solute-decorated population layer.
%
%  Standalone, compilable:  pdflatex -> bibtex -> pdflatex x2
%  Requires: amsmath, amssymb, booktabs, natbib[round,authoryear], hyperref
%  Bibliography: \bibliography{radcluster_2_1_RIS_plan}
% =====================================================================
\documentclass[11pt]{article}
\usepackage[margin=1in]{geometry}
\usepackage{amsmath,amssymb}
\usepackage{booktabs}
\usepackage[round,authoryear]{natbib}
\usepackage{xcolor}
\usepackage[colorlinks=true,linkcolor=blue,citecolor=blue,urlcolor=blue]{hyperref}
\usepackage{enumitem}

\newcommand{\cv}{c_{V}}
\newcommand{\ci}{c_{I}}
\newcommand{\Dv}{D_{V}}
\newcommand{\Di}{D_{I}}
\newcommand{\Cs}{C_{s}}
\newcommand{\code}[1]{\texttt{#1}}

\title{\textbf{RadCluster\_2\_1 Development Plan}\\[2pt]
\large Radiation-Induced Segregation and Solute Precipitation\\
as a Solute-Decorated Population Layer}
\author{RadCluster project --- planning note}
\date{June 2026}

\begin{document}
\maketitle

\begin{abstract}
This note plans the next minor version of RadCluster, \textbf{RadCluster\_2\_1},
which adds radiation-induced segregation (RIS) and irradiation-driven solute
precipitation to the existing cluster-dynamics framework for EUROFER97 /
ferritic--martensitic (F/M) steel. It (i) reviews a proposed design that
introduces RIS as a solute-decorated population/edge layer, (ii) synthesises
the RIS modelling literature and the data required, (iii) maps the model onto
the existing reaction-admissibility-graph (RAG) architecture, (iv) specifies a
tractable reduced model with conservation diagnostics, and (v) gives a phased
implementation plan with calibration and validation targets. The central
recommendation is endorsed --- RIS/precipitation should enter as graph
vertices and edges, not as an empirical hardening correction --- with three
substantive corrections: the RIS coupling coefficients must be grounded in
offline-tabulated Onsager/partial-diffusion transport ratios rather than free
fit parameters; the EUROFER97 calibration target must lead with Cr-rich
$\alpha'$ (Mn--Ni--Si precipitation is weak in EUROFER97 because Ni and Si are
deliberately minimised); and the C++ kernel-mirroring cost is the dominant
implementation risk and must be scheduled explicitly.
\end{abstract}

\tableofcontents

% =====================================================================
\section{Purpose and scope}
% =====================================================================
RadCluster\_2\_0 solves the coupled point-defect and cluster master equations
for SIA loops ($\tfrac12\langle111\rangle$ and $\langle100\rangle$), vacancy
cavities, and helium, with full-CD and bin-moment representations and two
helium-occupancy reductions \citep{Ghoniem2026}. Solutes (Cr, Mn, Si, W, C, N)
enter only as \emph{fixed} dissolved concentrations that scale the effective
jump frequencies $\omega^{\rm eff}_\alpha$ through the solute-trapping factor.

RadCluster\_2\_1 makes selected solutes \emph{dynamic} so that the model can
predict: (a) RIS --- the irradiation-driven enrichment/depletion of Cr, Mn,
Si, Ni at dislocations, loops, grain boundaries (GBs), and lath boundaries;
(b) the solute-rich nanofeatures that follow --- Cr-rich $\alpha'$ and
Mn--Ni--Si (MNS / G-phase) clusters; and (c) the precipitate contribution to
hardening. The design goal is to add this physics \emph{declaratively} ---
new populations, vertices, edges, and rate kernels --- while reusing the
existing solver, reductions, and conservation machinery.

% =====================================================================
\section{Review of the proposed design}
% =====================================================================
The proposed update introduces (1) bulk and sink-local solute reservoirs
$C_s^{b}$, $C_s^{\ell}$; (2) solute-rich precipitate populations $\alpha'$,
MNS (as a pseudo-species $G$), and loop-decorating clouds; (3) solute trapping,
emission, nucleation, radiation re-solution, and inter-population edges;
(4) per-solute conservation diagnostics; and (5) a dispersed-barrier hardening
sum. It correctly proposes to start from a reduced model rather than a full
$(n_{\rm Cr},n_{\rm Mn},n_{\rm Si},n_{\rm Ni})$ composition grid.

\subsection{What is correct and should be kept}
\begin{itemize}[leftmargin=1.4em]
\item \textbf{Architectural strategy.} Adding RIS/precipitation as
solute-decorated RAG vertices and edges is exactly how the framework is
designed to be extended. The edge taxonomy already contains
\code{SOLUTE\_TRAPPING}, \code{INTER\_POPULATION}, \code{SOURCE}, \code{SINK},
\code{COALESCENCE}, and \code{DISSOCIATION}; the state layout already supports
auxiliary scalar/array blocks (\code{add\_scalar}, \code{add\_aux}); and the
conservation block is already an auxiliary cumulative-integral pattern. No
solver change is required to add vertices and edges.
\item \textbf{Two-reservoir (bulk $+$ sink-local) RIS.} This is precisely the
0D-compatible reduction the literature recommends when a spatially resolved
profile across the sink is not solved \citep{Hu2021,Oplinger2026}. It reuses
the bulk $\cv,\ci$ and sink strengths RadCluster already computes each step.
\item \textbf{Pseudo-species $G$ for Mn--Ni--Si.} Collapsing the MNS
composition to a single coordinate avoids the multi-dimensional CD explosion
and is consistent with the fixed G-phase stoichiometry $\mathrm{Mn_6Ni_{16}Si_7}$
used in the benchmark CD studies \citep{Ke2018mns,Ke2018fm}.
\item \textbf{Re-solution term.} The athermal ejection form
$\Gamma_{\rm res}=b_0\,\dot\phi$ is the same mechanism already implemented as
P8 for helium and is the standard ballistic-mixing closure
\citep{NelsonHudsonMazey1972,Ke2019flux}.
\item \textbf{Per-solute conservation diagnostics and reduced-first
sequencing.} Direct analogues of $\delta_{\rm FP}$ and $\delta_{\rm He}$;
essential for debugging.
\end{itemize}

\subsection{Corrections required before implementation}
\label{sec:corrections}
\begin{enumerate}[leftmargin=1.6em]
\item \textbf{RIS coupling must come from transport coefficients, not free
$\eta$'s.} The proposed effective coefficients $\eta_{sI}^{\ell},
\eta_{sV}^{\ell}$ with hand-chosen signs hide the physics that sets both the
sign and the temperature dependence of RIS. The enrichment/depletion direction
is fixed by the \emph{ratio} of partial diffusion coefficients (equivalently
Onsager $L$-coefficients) for solute vs.\ matrix via the two mechanisms
\citep{WiedersichOkamotoLam1979,AllenWas1998,Messina2020,Oplinger2026}. These
must be tabulated offline as $d_{sV}(T),d_{sI}(T)$ (or the ratios) from the
\textsc{dft}{}+\textsc{scmf} dataset of \citet{Messina2014,Messina2020} and
\emph{read in}, not fitted. This also reproduces the Cr enrichment$\to$depletion
crossover automatically (Sec.~\ref{sec:ris-physics}).
\item \textbf{EUROFER97 calibration must lead with $\alpha'$, not MNS.} The
quoted target ($N\!\sim\!10^{24}\,\mathrm{m^{-3}}$, $d\!\sim\!3$--$5$~nm,
$32$~dpa, $332\,^\circ$C) corresponds to Mn--Ni--Si precipitation in
\emph{T91/RPV-type} steel \citep{Ke2018fm}. EUROFER97 deliberately minimises
Ni ($\lesssim0.01$~wt\%) and Si ($\sim0.05$~wt\%) \citep{Lindau2002}, so its
MNS driving force is weak. The dominant irradiation precipitate in high-Cr F/M
steel is Cr-rich $\alpha'$ \citep{Mathon2003,Bachhav2014,Reese2018}. The
primary EUROFER97 validation target should therefore be $\alpha'$
($N\!\sim\!10^{23}$--$10^{24}\,\mathrm{m^{-3}}$, $d\!\sim\!2$--$5$~nm); MNS
should be retained as a secondary, alloy-composition-gated channel and the
T91 MNS dataset used as a separate validation case.
\item \textbf{C++ kernel mirroring is the dominant cost.} The production RHS is
the hand-written C++ solver; it does \emph{not} walk the graph. Every new
population and edge must be mirrored by hand in
\code{rate\_kernels.cpp} for each \code{physics\_option} and each He case.
Adding RIS$+$precipitation roughly doubles the kernel surface. This is the main
schedule risk and is addressed in Sec.~\ref{sec:plan} by a Python-walker-first
strategy.
\item \textbf{Cr double role.} Cr is currently a fixed parameter inside the
solute-trapping factor that sets $\omega^{\rm eff}$. Once Cr is dynamic and
partitions into $\alpha'$, the bulk Cr that feeds trapping changes. For v2.1 we
recommend decoupling: keep the trapping factor on the (nearly constant) bulk Cr
to first order, and track a separate partitioned-Cr inventory; revisit full
self-consistency only if bulk Cr depletion exceeds a few percent.
\item \textbf{General \code{SOLUTE\_TRAPPING} walker handler.} In the current
graph walker the \code{SOLUTE\_TRAPPING} handler raises
\code{NotImplementedError} (He is handled by its dedicated reduction). A
general composition-changing handler (or a solute reduction analogous to the He
Case-1 mean-field) must be implemented; this is real work the design understates.
\end{enumerate}

% =====================================================================
\section{RIS physics background}
\label{sec:ris-physics}
% =====================================================================
\subsection{Mechanisms and direction}
RIS arises because irradiation drives persistent vacancy and interstitial
fluxes to sinks; solutes couple to these fluxes
\citep{WiedersichOkamotoLam1979,Marwick1978,Was2017,ArdellBellon2016}. Two mechanisms
compete: the \emph{inverse-Kirkendall} (vacancy) mechanism, in which solutes
that exchange faster with vacancies deplete at sinks; and the
\emph{interstitial-association / drag} mechanism, in which solutes carried in
mobile mixed dumbbells (or dragged by vacancies) enrich at sinks. In bcc Fe,
\textsc{dft}+\textsc{scmf} transport theory predicts that P, Mn, and Cr diffuse
preferentially as dumbbells (enrichment), while Cu, Ni, Si enrich mainly by
vacancy drag \citep{Messina2020}; the fcc Ni-base austenitic surrogate behaves
analogously but with several sign reversals \citep{Toijer2021}, and atomistic
kinetic Monte Carlo for Fe--Cr reproduces the radiation-accelerated Cr diffusion
and $\alpha'$ formation that follow \citep{Senninger2016,SoissonJourdan2016}.
Cr is a fine balance with a crossover:
\emph{enrichment below $\sim$543~K, depletion above}
\citep{Messina2020}; experiments on F/M steel show Cr enrichment by the
interstitial mechanism over a wide range with depletion only at the highest
temperatures \citep{WharryWas2014,WharryWas2013}. RIS magnitude is bell-shaped
in temperature and shifts to higher $T$ with dose rate; temperature largely
sets the \emph{direction}, while dose rate, sink density, and grain size set the
\emph{magnitude} \citep{Oplinger2026}.

\subsection{Governing equations (the implementable core)}
For a sink-normal coordinate the coupled balances are
\citep{WiedersichOkamotoLam1979,AllenWas1998}
\begin{align}
\frac{\partial \Cs}{\partial t} &= -\nabla\!\cdot J_s, \\
\frac{\partial \cv}{\partial t} &= -\nabla\!\cdot J_V + K_0 - R\,\cv\ci - k_V^2 \Dv(\cv-\cv^{\rm eq}),\\
\frac{\partial \ci}{\partial t} &= -\nabla\!\cdot J_I + K_0 - R\,\cv\ci - k_I^2 \Di\,\ci,
\end{align}
with recombination $R=4\pi r_0(\Dv+\Di)/\Omega$ and sink strengths $k^2$. The
solute flux that produces segregation is the difference of the vacancy and
interstitial contributions,
\begin{equation}
\boxed{\,
J_s = -\big(d_{sV}+d_{sI}\big)\nabla \Cs
      \;-\;\Cs\big(d_{sV}\,\nabla\!\ln\cv \;-\; d_{sI}\,\nabla\!\ln\ci\big),
\,}
\label{eq:Js}
\end{equation}
where the partial diffusion coefficients factorise as
$d_{sd}=\alpha_{sd}(T)\,c_d$ with
$\alpha_{sd}=f_d\lambda^2 z\,\nu_{sd}\exp(-E_m^{sd}/k_BT)$
(correlation factor $f_d$, jump distance $\lambda$, coordination $z$).
The \emph{sign} of the bracket in Eq.~\eqref{eq:Js} --- set by the ratio
$d_{sI}/d_{\rm Fe,I}$ vs.\ $d_{sV}/d_{\rm Fe,V}$ --- determines enrichment vs.\
depletion. Modern practice replaces the empirical $d_{sd}$ by Onsager
$L$-coefficients computed offline by \textsc{scmf}/\textsc{KineCluE}
\citep{NastarSoisson2020,SchulerMessinaNastar2020}; the rate-theory form of
\citet{Oplinger2026} writes the same physics with relative transport
coefficients and shows two cascade biases that matter for a fusion/fission
code: a \emph{production bias} (excess freely migrating vacancies because SIAs
cluster into glissile loops --- already modelled here) and an
\emph{absorption bias} ($Z_I>Z_V$).

\subsection{0D reduction: the sink-local reservoir}
\label{sec:0d-reduction}
A pure 0D code cannot resolve the $\sim$5--20~nm GB profile. The tractable
reduction lumps each sink class $\ell$ into a reservoir whose solute fraction
relaxes under the net RIS current driven by the \emph{bulk} defect
supersaturations RadCluster already computes:
\begin{equation}
\boxed{\;
\frac{dC_s^{\ell}}{dt}
= \frac{S_\ell}{w_\ell}\Big[
   \underbrace{-\,d_{sV}\,\Dv(\cv-\cv^{\rm eq})}_{\text{IK back-flux}}
   \;+\;\underbrace{d_{sI}\,\Di\,\ci}_{\text{drag}}
   \;-\;\underbrace{k^{\rm back}_{s}\big(C_s^{\ell}-C_s^{b}\big)}_{\text{thermal back-diffusion}}
  \Big],
\;}
\label{eq:reservoir}
\end{equation}
with $S_\ell$ the sink strength of class $\ell$, $w_\ell$ an effective
segregation width ($\sim$5--10~nm), and $k^{\rm back}_s$ a relaxation rate. The
two flux terms carry opposite signs for the two mechanisms; the back-diffusion
term gives the high-$T$ roll-off. Eq.~\eqref{eq:reservoir} naturally produces
the bell-shaped temperature response without a spatial mesh
\citep{Hu2021,Oplinger2026}. The 1D PDE (MIK, Oplinger) is retained
\emph{offline} only to calibrate $w_\ell$ and $k^{\rm back}_s$.

\subsection{RIS-driven precipitation and re-solution}
Local enrichment from Eq.~\eqref{eq:reservoir} raises the sink-local
supersaturation above the bulk solubility and triggers heterogeneous
nucleation that does not occur in the bulk. The benchmark MNS CD study shows
this explicitly: without RIS, MNS clusters cannot grow past $\sim$70 atoms;
RIS supplies the supersaturation for nucleation and growth on dislocations
\citep{Ke2018mns,Ke2018fm}; atom-probe tomography confirms solute clouds
decorating dislocation loops as the heterogeneous nucleation sites
\citep{Davis2020}. The counter-term is ballistic re-solution: cascades
eject solute back to the matrix, capping the size. We adopt the cluster-dynamics
form already present for He,
\begin{equation}
\left.\frac{dC_n^{p}}{dt}\right|_{\rm res} = -\,b_0\,\dot\phi\,n\,C_n^{p}
\quad(\text{gain at the monomer size}),
\label{eq:resolution}
\end{equation}
with the modern calibration $b\!\approx\!50$ replacements per dpa over a
relocation distance $R\!\approx\!0.25$~nm for solute precipitates
\citep{Ke2019flux,SchwenAverback2016}.

% =====================================================================
\section{Required data}
\label{sec:data}
% =====================================================================
Tables~\ref{tab:transport}--\ref{tab:thermo} collect the parameters needed for
the reduced model, with primary sources. Values are dilute-limit
\textsc{dft}+\textsc{scmf} unless noted; alloy-concentration corrections shift
the Cr crossover and should be applied where available.

\begin{table}[h]
\centering\small
\caption{bcc-Fe solute transport / RIS direction \citep{Messina2014,Messina2020}.
$E_m$ are migration barriers; ``RIS'' is the predicted sink behaviour with
crossover temperature where applicable.}
\label{tab:transport}
\begin{tabular}{lccccl}
\toprule
Solute & $E_m^{\rm vac}$ (eV) & $E_m^{\rm db}$ (eV) & $F_B^{\rm V}$ (eV) &
  dominant mech. & RIS at sinks \\
\midrule
Cr & 0.65 & 0.22 & 0.124 & dumbbell/vac (balance) & enrich $<\!543$~K; deplete above \\
Mn & 0.62 & 0.32 & 0.235 & dumbbell & enrich \\
Si & 0.72 & 0.82 & 0.223 & vacancy drag & enrich \\
Ni & 0.70 & 0.67 & $-0.008$ & vacancy drag & enrich ($<\!1089$~K) \\
P  & 0.67 & 0.21 & 0.276 & dumbbell & enrich \\
Cu & 0.69 & 0.88 & 0.205 & vacancy & enrich ($<\!1086$~K) \\
\bottomrule
\end{tabular}
\end{table}

\begin{table}[h]
\centering\small
\caption{Measured irradiation precipitate populations (calibration/validation
targets).}
\label{tab:ppt}
\begin{tabular}{llllll}
\toprule
Phase & Alloy & $T$ & dose & $N$ (m$^{-3}$) & $d$ (nm) \\
\midrule
$\alpha'$ & Fe-15Cr & 320$^\circ$C & 7 dpa & $1.1\times10^{24}$ & $\sim$2.6 \citep{Reese2018}\\
$\alpha'$ & Fe-9Cr  & 320$^\circ$C & 7 dpa & $1.2\times10^{23}$ & $\sim$4.0 \citep{Reese2018}\\
$\alpha'$ & Fe-18Cr & 290$^\circ$C & 1.8 dpa & $\sim5\times10^{24}$ & small \citep{Bachhav2014}\\
MNS/G    & T91     & 400$^\circ$C & $\sim$7 dpa & $1.3\times10^{23}$ & $\sim$3.2 \citep{Ke2018mns}\\
MNS/G    & T91 (n) & $\sim$332$^\circ$C & $\sim$32 dpa & $\sim10^{24}$ & 3--5 \citep{Ke2018fm}\\
\bottomrule
\end{tabular}
\end{table}

\begin{table}[h]
\centering\small
\caption{Thermodynamic, interfacial, and re-solution data.}
\label{tab:thermo}
\begin{tabular}{lll}
\toprule
Quantity & Value & Source \\
\midrule
Cr solubility in $\alpha$-Fe (300$^\circ$C) & $\sim$9--10 at\% & \citep{Bonny2008,Mathon2003}\\
$\alpha'$ composition & 80--96 at\% Cr & \citep{Reese2018}\\
$\alpha'$ interfacial energy & 0.090 J\,m$^{-2}$ & \citep{Ke2019flux}\\
G-phase interfacial energy & 0.19 J\,m$^{-2}$ & \citep{Ke2018mns}\\
G-phase stoichiometry & $\mathrm{Mn_6Ni_{16}Si_7}$ (L2$_1$) & \citep{Ke2018mns}\\
Ballistic mixing $b$ & 50 replacements/dpa & \citep{Ke2019flux}\\
Relocation distance $R$ & 0.249 nm & \citep{Ke2019flux}\\
Cascade survival $\xi$ & 0.33 (neutron/ion) & \citep{Ke2019flux}\\
\bottomrule
\end{tabular}
\end{table}

\noindent\textbf{EUROFER97 reservoir initial conditions} (nominal, Fe balance)
\citep{Lindau2002}: Cr $\approx$ 8.9 wt\% ($\sim$9.5 at\%), Mn $\approx$ 0.42
wt\%, W $\approx$ 1.1 wt\%, V $\approx$ 0.19 wt\%, Si $\approx$ 0.05 wt\%,
Ni $\lesssim$ 0.01 wt\%, C $\approx$ 0.11 wt\%.

\paragraph{Data gaps to close before calibration.}
(i) The exact $N,d$ pair for the $32$~dpa / $332\,^\circ$C MNS case is in the
paywalled \citet{Ke2018fm}; verify from the primary ATR-2 APT table.
(ii) A self-consistent Fe--Cr regular-solution $\Omega$ / CALPHAD parameter set
(Andersson--Sundman / SGTE) for the $\alpha'$ driving force.
(iii) Mn engineering diffusivity ($D_0,E_a$) --- use \citet{Messina2020}
Table values rather than the MNS-paper SI, which omits Mn.
(iv) The DFT Cr crossover ($\sim$543~K) is below the experimental Cr-enrichment
regime ($\lesssim$600$^\circ$C); apply alloy-concentration/SIA-flux corrections,
not the dilute value alone.

% =====================================================================
\section{Integration into the RAG architecture}
\label{sec:arch}
% =====================================================================
The model is declared in \code{py\_utils/materials/eurofer97/declaration.py}
through the existing API; no change to \code{rag.py}, \code{state\_layout.py},
or the integrator is needed to add vertices/edges.

\subsection{New populations and reservoirs}
\begin{itemize}[leftmargin=1.4em]
\item Extend the decoration list:
\code{ReactionAdmissibilityGraph("EUROFER-97", gas\_species=["He","Cr","Mn","Si","Ni"])}.
\item Precipitate populations via \code{rag.add\_population(Population(...))}:
\code{alpha\_prime} (vacancy-type, immobile, $n_{\min}\!\gtrsim\!2$),
\code{MNS} (pseudo-species $G$, immobile), and \code{loop\_cloud}
(decoration of the $\tfrac12\langle111\rangle$ population).
\item Solute reservoirs as scalar aux blocks via \code{layout.add\_scalar}:
one bulk reservoir $C_s^{b}$ per dynamic solute, and one sink-local reservoir
$C_s^{\ell}$ per (solute, sink-class) pair to be tracked. Start with
$\ell\in\{\rm disl,loop,GB\}$.
\end{itemize}

\subsection{New edge classes and kernels}
All map onto existing \code{EdgeClass} values; kernels are precomputed arrays
registered with \code{rag.register\_kernel}:
\begin{itemize}[leftmargin=1.4em]
\item \code{SOURCE}/\code{SINK} pair on the reservoirs implementing the RIS
exchange of Eq.~\eqref{eq:reservoir} (``\code{RIS\_TRANSFER}'').
\item \code{SOLUTE\_TRAPPING} (trap direction): precipitate growth by solute
capture, $\beta_{n,s}=4\pi r_n Z\,D_s^{\rm irr}$ with
$D_s^{\rm irr}=D_s^{\rm th}+D_s^{\rm RED}(\cv,\ci)$.
\item \code{SOLUTE\_TRAPPING} (detrap direction): thermal emission set by
detailed balance to the correct solubility; plus athermal re-solution
Eq.~\eqref{eq:resolution}.
\item \code{SOURCE} for nucleation ($\varnothing\!\to\!P_{n_*}$), rate gated by
the sink-local supersaturation $(C_s^{\ell}-C_s^{\rm eq})_+$.
\item \code{COALESCENCE} for precipitate coarsening and
\code{INTER\_POPULATION}/\code{LOOP\_DECORATION} for loop-cloud formation
proportional to loop density.
\end{itemize}

\subsection{State layout and reductions}
Precipitate populations use the existing \code{add\_discrete} or
\code{add\_bin\_moment} blocks (the bin-moment reduction applies unchanged ---
solute clusters span decades in size). Reservoirs are scalar aux blocks. A new
\code{conservation} aux entry per dynamic solute holds the cumulative
source/loss integral. The required new \emph{reduction} work is a general
\code{SOLUTE\_TRAPPING} graph-walker handler (Sec.~\ref{sec:corrections},
item~5), modelled on the He Case-1 mean-field loading.

\subsection{C++ mirror}
Each new edge must be added to \code{cpp\_utils/materials/eurofer97/rate\_kernels.cpp}
and the parameters serialised in \code{cpp\_bridge.py}. Because this is the bulk
of the effort, v2.1 develops and validates the physics through the Python graph
walker first (Sec.~\ref{sec:plan}).

% =====================================================================
\section{The RadCluster\_2\_1 reduced model}
\label{sec:model}
% =====================================================================
\paragraph{State (minimal first cut).}
Dynamic solutes $s\in\{\mathrm{Cr},\mathrm{Mn},\mathrm{Si},\mathrm{Ni}\}$ with
bulk reservoirs $C_s^{b}$ and sink-local reservoirs $C_s^{\ell}$,
$\ell\in\{\rm disl,loop,GB\}$; precipitate size distributions $C_n^{\alpha'}$,
$C_n^{G}$, $C_n^{\rm cloud}$. This adds $\mathcal{O}(4{+}12)$ scalar reservoir
ODEs plus three cluster ladders (each bin-moment-reduced) to the existing
system.

\paragraph{Governing equations.}
Reservoirs evolve by Eq.~\eqref{eq:reservoir} (sink-local) and its bulk
complement; precipitate distributions follow the standard RadCluster master
form
\begin{equation}
\frac{dC_n^{p}}{dt}=G_n^{p}
 +\sum_s\!\big[\beta_{n-1,s}^{p}C_{n-1}^{p}C_s^{\ell}-\beta_{n,s}^{p}C_n^{p}C_s^{\ell}\big]
 +\sum_s\!\big[\alpha_{n+1,s}^{p}C_{n+1}^{p}-\alpha_{n,s}^{p}C_n^{p}\big]
 -\Gamma_{\rm res}^{p}C_n^{p}-D_n^{p}C_n^{p},
\end{equation}
which is exactly $\dot c_a=G_a+\sum_r S_{ar}J_r-D_a c_a$ --- no solver change.
$D_s^{\rm RED}$ in $\beta$ is tied to the $\cv,\ci$ already integrated.

\paragraph{Coupling back to defects and hardening.}
Precipitates add to the sink strength,
$k_{\rm ppt}^2=\sum_{p,n}4\pi r_n^{p}Z_p\,C_n^{p}$, fed into the existing
point-defect balance; and to hardening through a dispersed-barrier sum
$\Delta\sigma_{\rm ppt}=M\mu b\,[\sum_{p,k}(\alpha_{p,k}^2 N_{p,k}d_{p,k})]^{1/2}$,
combined in quadrature with the existing loop and cavity contributions.

% =====================================================================
\section{Conservation diagnostics}
% =====================================================================
For each dynamic solute $s$ define the total inventory
\begin{equation}
C_s^{\rm tot}=C_s^{b}+\sum_\ell C_s^{\ell}
  +\sum_{p,n} n_s(n)\,C_n^{p},
\end{equation}
and require $dC_s^{\rm tot}/dt=0$ (no transmutation source/loss for these
solutes). The diagnostic
$\delta_s=|C_s^{\rm tot}(t)-C_s^{\rm tot}(0)|/C_s^{\rm tot}(0)$ is the direct
analogue of $\delta_{\rm FP}$ and $\delta_{\rm He}$ and must stay
$\lesssim10^{-6}$. As with helium, each sink/exchange flux removed from one
reservoir must be added to its partner so the ledger closes by construction;
the cumulative integrals live in the \code{conservation} aux block. This is the
same discipline that exposed and fixed the helium balance bugs in v2.0 and is a
hard acceptance gate for v2.1.

% =====================================================================
\section{Phased implementation plan}
\label{sec:plan}
% =====================================================================
\begin{description}[leftmargin=0em,style=nextline]
\item[Phase 0 --- Prerequisites.] Resolve the pre-existing $\mathcal{O}(I^2)$
\code{phi\_junc} allocation so production sizes run; assemble the offline
transport table $d_{sV}(T),d_{sI}(T)$ from \citet{Messina2014,Messina2020} and
the $\alpha'$/G-phase thermodynamic set (Sec.~\ref{sec:data} gaps).
\item[Phase 1 --- Dynamic solute reservoirs + RIS, Python walker.] Add bulk and
sink-local reservoirs and the \code{RIS\_TRANSFER} edges
(Eq.~\eqref{eq:reservoir}). Implement the general \code{SOLUTE\_TRAPPING}
handler. \emph{Acceptance:} per-solute $\delta_s\!<\!10^{-6}$; reproduce the
bell-shaped $T$ dependence and the Cr enrichment$\to$depletion crossover.
\item[Phase 2 --- $\alpha'$ precipitation.] Add the \code{alpha\_prime}
population with capture/emission (detailed balance to Fe--Cr solubility),
nucleation gated by sink-local Cr, and re-solution Eq.~\eqref{eq:resolution}.
\emph{Acceptance:} Cr conservation holds; calibrate to the Fe--(9--18)Cr
$\alpha'$ data (Table~\ref{tab:ppt}).
\item[Phase 3 --- MNS pseudo-species and loop clouds.] Add the $G$ population
and \code{LOOP\_DECORATION}; gate strength on EUROFER97's low Ni/Si.
\emph{Acceptance:} weak MNS in EUROFER97; quantitative match to the T91 MNS case
as a separate validation.
\item[Phase 4 --- Defect/hardening coupling.] Feed $k_{\rm ppt}^2$ into the
point-defect balance and add $\Delta\sigma_{\rm ppt}$.
\item[Phase 5 --- C++ mirror and performance.] Port the validated edges to
\code{rate\_kernels.cpp} and \code{cpp\_bridge.py}; bit-for-bit cross-check
against the Python walker on a small system before production runs.
\end{description}

\paragraph{Validation targets.} Cr-rich $\alpha'$ in Fe--(9--18)Cr
\citep{Reese2018,Bachhav2014}; GB RIS magnitudes and the Cr crossover in F/M
steel \citep{WharryWas2014}; and the T91 MNS case \citep{Ke2018fm}. Each phase
gates on its conservation diagnostic.

% =====================================================================
\section{Risks and open decisions}
% =====================================================================
\begin{itemize}[leftmargin=1.4em]
\item \textbf{C++ mirroring effort} (Sec.~\ref{sec:corrections}, item~3) ---
mitigated by Python-walker-first development; decision needed on whether v2.1
ships the Python walker as a supported solver path or treats it as reference
only.
\item \textbf{RIS reduction fidelity} --- the lumped reservoir
(Eq.~\eqref{eq:reservoir}) cannot capture the GB profile width; calibrate
$w_\ell,k^{\rm back}_s$ offline against a 1D MIK/Oplinger solve.
\item \textbf{Cr self-consistency} (Sec.~\ref{sec:corrections}, item~4) ---
decide the threshold bulk-Cr depletion at which the trapping factor must be
recomputed self-consistently.
\item \textbf{MNS in EUROFER97} --- confirm the modelled heat composition;
EUROFER97's low Ni/Si may make MNS negligible, in which case Phase~3 is
validation-only against T91.
\item \textbf{Data provenance} --- close the four data gaps in
Sec.~\ref{sec:data} before Phase~2 calibration.
\end{itemize}

\bibliographystyle{plainnat}
\bibliography{radcluster_2_1_RIS_plan}

\end{document}
